% Options for packages loaded elsewhere
\PassOptionsToPackage{unicode}{hyperref}
\PassOptionsToPackage{hyphens}{url}
%
\documentclass[
  12pt,
]{article}
\usepackage{amsmath,amssymb}
\usepackage{setspace}
\usepackage{iftex}
\ifPDFTeX
  \usepackage[T1]{fontenc}
  \usepackage[utf8]{inputenc}
  \usepackage{textcomp} % provide euro and other symbols
\else % if luatex or xetex
  \usepackage{unicode-math} % this also loads fontspec
  \defaultfontfeatures{Scale=MatchLowercase}
  \defaultfontfeatures[\rmfamily]{Ligatures=TeX,Scale=1}
\fi
\usepackage{lmodern}
\ifPDFTeX\else
  % xetex/luatex font selection
  \setmainfont[]{TeX Gyre Pagella}
\fi
% Use upquote if available, for straight quotes in verbatim environments
\IfFileExists{upquote.sty}{\usepackage{upquote}}{}
\IfFileExists{microtype.sty}{% use microtype if available
  \usepackage[]{microtype}
  \UseMicrotypeSet[protrusion]{basicmath} % disable protrusion for tt fonts
}{}
\makeatletter
\@ifundefined{KOMAClassName}{% if non-KOMA class
  \IfFileExists{parskip.sty}{%
    \usepackage{parskip}
  }{% else
    \setlength{\parindent}{0pt}
    \setlength{\parskip}{6pt plus 2pt minus 1pt}}
}{% if KOMA class
  \KOMAoptions{parskip=half}}
\makeatother
\usepackage{xcolor}
\usepackage{graphicx}
\usepackage[margin=1in]{geometry}
\usepackage{longtable,booktabs,array}
\usepackage{calc} % for calculating minipage widths
% Correct order of tables after \paragraph or \subparagraph
\usepackage{etoolbox}
\makeatletter
\patchcmd\longtable{\par}{\if@noskipsec\mbox{}\fi\par}{}{}
\makeatother
% Allow footnotes in longtable head/foot
\IfFileExists{footnotehyper.sty}{\usepackage{footnotehyper}}{\usepackage{footnote}}
\makesavenoteenv{longtable}
\setlength{\emergencystretch}{3em} % prevent overfull lines
\providecommand{\tightlist}{%
  \setlength{\itemsep}{0pt}\setlength{\parskip}{0pt}}
\setcounter{secnumdepth}{-\maxdimen} % remove section numbering
\usepackage[natbibapa]{apacite}
\usepackage{lineno}\linenumbers
\ifLuaTeX
  \usepackage{selnolig}  % disable illegal ligatures
\fi
\IfFileExists{bookmark.sty}{\usepackage{bookmark}}{\usepackage{hyperref}}
\IfFileExists{xurl.sty}{\usepackage{xurl}}{} % add URL line breaks if available
\urlstyle{same}
\hypersetup{
  pdftitle={Cross-Cultural Validation and Ideological Fairness of a Historical Perspective Taking Instrument: Evidence from Czech Secondary Students},
  pdfauthor={Juda Kaleta},
  hidelinks,
  pdfcreator={LaTeX via pandoc}}

\title{Cross-Cultural Validation and Ideological Fairness of a Historical Perspective Taking Instrument: Evidence from Czech Secondary Students}
\author{Juda Kaleta}
\date{2026}

\begin{document}
\maketitle

\setstretch{2}
Juda Kaleta

Institute of History, Faculty of Arts, Charles University, Czech Republic

Email: juda.kaleta@ff.cuni.cz | ORCID: 0000-0003-0043-4039

\hypertarget{abstract}{%
\subsection{Abstract}\label{abstract}}

This study reports the first Czech adaptation of the Hartmann and
Hasselhorn Historical Perspective Taking (HPT) instrument and examines
whether its politically charged scenario introduces
construct-irrelevant variance (CIV) related to students' ideological
attitudes. Among 293 Czech secondary students, confirmatory factor
analysis favored a correlated three-factor representation. An item-level
bias analysis detected no differential item functioning across ideology
groups with this sample and procedure. Multi-group factor analysis
produced fit-index changes compatible with scalar invariance, but
inadmissible baseline solutions prevent a firm equivalence conclusion. Historical knowledge
correlated with the primary HPT\_CTX6 composite (\emph{r} = .34, 95\% CI [.23, .43]); ideology was
unrelated (\emph{r} = $-$.05, 95\% CI [$-$.16, .07]). These findings support a
possible protocol for evaluating measurement fairness when assessment
content carries attitudinal valence, subject to the limits of this
single scenario, item format, and ideologically restricted sample.

\textbf{Keywords:} historical perspective taking, cross-cultural validation, measurement invariance, differential item functioning, construct-irrelevant variance, test fairness

\hypertarget{introduction}{%
\subsection{1. Introduction}\label{introduction}}

Cross-cultural adaptation of educational and psychological instruments
requires more than linguistic translation: it demands evidence that
the adapted version measures the same construct, with equivalent
psychometric properties, in the target population \citep{AERA2014Standards,Hambleton2005,ITC2017}. When assessment content carries
attitudinal valence---when scenarios, stimuli, or items touch on
beliefs that respondents may hold---an additional fairness question
arises: does the content introduce construct-irrelevant variance (CIV)
by functioning differently for subgroups whose attitudes align with
that content \citep{Messick1995Validity}? This question is especially pressing for
scenario-based instruments in the social sciences, where rich
contextual material is needed to elicit complex reasoning but may
simultaneously activate respondents' pre-existing attitudes.
Historical perspective taking (HPT)---the capacity to reconstruct the
beliefs, values, and circumstances that made past actors' decisions
intelligible in their own time \citep{Hartmann_Hasselhorn2008Historical,Lee_Ashby2001Empathy}---provides a particularly informative case. HPT
instruments routinely use politically and morally charged historical
scenarios as the stimulus for disciplinary reasoning (reasoning
according to the methodological norms of the historical discipline,
such as source criticism and contextualization; \citealp{VanDrie_VanBoxtel2008Historical}; \citealp{Seixas_Morton2013TheBig}; \citealp{Wineburg2001Historical}).

The instrument reported by \citet{Hartmann_Hasselhorn2008Historical} provides a
scenario-based measure of HPT. It was validated with German Grade 10
students \citep{Hartmann_Hasselhorn2008Historical} and subsequently
translated and studied with Dutch students aged 10--17
\citep{Huijgen_et_al2014Testing,Huijgen_et_al2017Toward}. This evidence
base is promising but remains concentrated in two language contexts
and a small number of studies. The instrument operationalizes HPT through a scenario-based
approach in which students read about a historically situated dilemma
and respond to nine Likert-scaled items in three interpretive
response categories: present-oriented perspective (POP), role of the
historical agent (ROA), and historical contextualization (CONT). POP
and CONT represent the clearest contrast between present-centered
judgment and contextual reasoning; ROA's position as an intermediate
level remains theoretically less certain
\citep{Hartmann_Hasselhorn2008Historical,Lee_Ashby2001Empathy}.

The present study reports the first Czech adaptation and psychometric evaluation of the Hartmann and Hasselhorn HPT instrument. Two considerations motivate this work. First, no validated Czech-language HPT instrument was identified. Czech initiatives have developed digital historical-literacy tasks and measures of epistemic beliefs about history \citep{Cinatl_et_al2021Historicka,Ripka_et_al2024Epistemic}; the present adaptation provides a candidate instrument for further Czech validation and future cross-national equivalence studies.

Second, the instrument's scenario content raises a specific fairness question. Hannes discusses Germany's crisis and the 1930 election; students judge how well nine explanations fit his situation, including whether he could vote for an anti-democratic party such as the NSDAP. These judgments could, in principle, function differently for students whose political attitudes align with the scenario's ideological content. Students with right-authoritarian sympathies might find it cognitively easier to endorse contextualizing responses---not because they are reasoning historically, but because the responses resonate with their own attitudes. If present, such a pattern would constitute construct-irrelevant variance (CIV): systematic score variation attributable to a factor extraneous to the target construct \citep{AERA2014Standards,Messick1995Validity}.

This concern is not merely theoretical. Research on identity-protective cognition demonstrates that individuals process information through the lens of pre-existing beliefs \citep{Jost2003,Kahan2017}. Recent work has documented how ideological content can infiltrate competence judgments in history assessment, whether through teacher scoring bias in open-ended narrative
assessment \citep{Alven2019Bias} or through construct comparability challenges in cross-cultural settings \citep{Badham_et_al2025Construct}. The question central to the present study is whether the HPT instrument's psychometric properties---its factor structure, item parameters, and score interpretations---hold equivalently across students with different political attitudes, or whether politically charged content introduces measurement bias.

The study contributes to the literature in three ways: (a) it applies a possible methodological approach for evaluating whether attitudinally charged assessment content introduces construct-irrelevant variance, (b) it provides DIF and measurement-invariance evidence across political-ideology groups, and (c) it adds Czech evidence to the small cross-cultural evidence base for this measure. To establish this evidence, we conducted confirmatory factor analysis (CFA), reliability analysis, intraclass correlation estimation, IRT-based differential item functioning (DIF) analysis, and multi-group CFA for measurement invariance testing.

\hypertarget{literature-review}{%
\subsection{2. Literature Review}\label{literature-review}}

\hypertarget{historical-perspective-taking-and-its-measurement}{%
\subsubsection{2.1 Historical Perspective Taking and Its Measurement}\label{historical-perspective-taking-and-its-measurement}}

HPT refers to the ability to understand past actors' decisions by
reconstructing the historical context in which they operated \citep{Hartmann_Hasselhorn2008Historical,Lee_Ashby2001Empathy}. HPT is a second-order
(procedural or disciplinary) concept---as distinct from first-order
(substantive) concepts such as ``revolution'' or ``empire,'' it
concerns the meta-level reasoning processes historians use to make
sense of the past \citep{Seixas_Peck2004Teaching}. The concept has roots in
British traditions that progressively refined the construct from
historical empathy \citep{Ashby1987,Shemilt1984} toward
rational reconstruction \citep{Lee_Ashby2001Empathy}. It has been situated
within broader frameworks of historical reasoning that encompass
evidence use, argumentation, and substantive concept deployment---the
application of domain-specific content knowledge about events, causes,
and actors \citep{vanBoxtel_vanDrie2018Historical,VanDrie_VanBoxtel2008Historical,VanSledright2014Assessing}. The literature variously labels this
capacity ``historical empathy,'' ``perspective taking,'' or ``rational
reconstruction.'' We follow \citet{Hartmann_Hasselhorn2008Historical} in adopting
the term ``perspective taking'' within the rationalist tradition of
\citet{Lee_Ashby2001Empathy}, which emphasizes cognitive reconstruction of
context rather than affective identification with past actors (see
\citealp{Barton2004}, for a contrasting view that emphasises
affective engagement alongside cognitive reconstruction; see also
\citealp{Endacott_Brooks2013An_Updated,Endacott2014Negotiating}, for
treatments of the empathy/perspective-taking distinction). The present
instrument therefore operationalises cognitive perspective taking, not
the fuller cognitive-affective construct of historical empathy. Across these frameworks, HPT
is understood as a cognitively demanding, knowledge-dependent
competence---not a disposition but a disciplinary achievement requiring
both domain-specific knowledge and the capacity to deploy that
knowledge in reconstructing past perspectives \citep{Chapman2021Introduction,VanSledright2014Assessing,Wineburg2001Historical}. Distinguishing genuine
disciplinary reasoning from attitude-driven responding is a
measurement challenge that extends beyond HPT to epistemic cognition in
history more broadly \citep{Maggioni_et_al2009Walking}.

\citet{Hartmann_Hasselhorn2008Historical} operationalized HPT using three response
categories: \emph{present-oriented perspective} (POP), where students judge
historical actors by contemporary standards; \emph{role of the
historical agent} (ROA), where students explain action through familiar
or stereotyped social roles; and
\emph{contextualization} (CONT), where students situate decisions
within the specific historical conditions that rendered them
intelligible. The POP--CONT continuum from present-centered judgment
toward contextualization has empirical parallels in earlier
work on students' historical understanding \citep{Lee_Shemilt2003A_scaffold}. The
original published validation study (\emph{n} = 170 German students,
Grade 10; \citealp{Hartmann_Hasselhorn2008Historical})
placed raw POP and CONT at opposing poles of one principal component
and ROA on a second, less clearly interpreted component (explained
variance = 51\%); a separate latent-class analysis identified three
student classes. Students classified in the highest class had
significantly better history grades; however, PCA is an exploratory
technique, and the present study tests a correlated three-factor
representation rather than reproducing the original PCA dimensions.
\citet{Huijgen_et_al2014Testing} replicated the two-dimensional pattern
for the Nazi-party scenario with Dutch students, while
\citet{Huijgen_et_al2017Toward} combined item ratings and think-alouds
to examine how students interpreted the task.

The difficulty of measuring HPT stems from its inherently interpretive character. Unlike factual recall, HPT cannot be adequately inferred from decontextualized factual-recall items alone \citep{Breakstone2014Try,Smith2017Cognitive}. Scenario-based instruments address this challenge by presenting contextual information and requiring students to evaluate statements reflecting different levels of historical reasoning \citep{Hartmann_Hasselhorn2008Historical,Seixas_et_al2015A_Design}. However, this approach introduces a measurement tension: the richer the historical context, the more opportunity for construct-irrelevant factors---including attitudes toward that content---to influence responses.

\hypertarget{cross-cultural-test-adaptation}{%
\subsubsection{2.2 Cross-Cultural Test Adaptation}\label{cross-cultural-test-adaptation}}

The adaptation of educational and psychological instruments across languages and cultures requires systematic procedures to ensure that the adapted version measures the same construct as the original \citep{Hambleton2005,ITC2017}. The International Test Commission (ITC) guidelines emphasize that adaptation involves more than translation: it requires attention to construct equivalence, cultural appropriateness, and technical quality of the adapted version \citep{ITC2017}. For instruments used in history education, cross-cultural adaptation faces the additional challenge that historical content may carry different resonances across national contexts. A scenario about the Weimar Republic, originally designed for German
students, intersects differently with Czech historical memory---while
subsequent events shape contemporary Czech memory of the interwar
period, the scenario centres on Weimar-era deliberation (early 1930s),
which predates the 1938 Munich Agreement and cession of the
Sudetenland by approximately eight years. Nevertheless, the broader
period connects directly to the Nazi occupation of 1939--1945 in Czech
collective memory.

In Czech lower-secondary education, history is prescribed within the
compulsory People and Society curriculum area for Grades 6--9. The
framework includes the 1920s and 1930s, fascism, Nazism, totalitarianism,
and work with historical sources, but does not define HPT as a separate
competence \citep[pp.~54, 57]{RVPZV2023}. School curricula broadly align
with this framework, although curriculum documents describe intended
rather than enacted provision \citep{Kaleta2025Sticking}. Participating
teachers reported that students had encountered the relevant interwar
content; this does not establish equivalent prior instruction across
schools.

Cross-cultural validation studies should demonstrate that the adapted instrument preserves the factor structure, reliability, and validity relationships of the original \citep{Hambleton2005}. At minimum, this requires confirming the theorized factor structure through CFA, reporting reliability estimates appropriate for the item characteristics, and providing convergent validity evidence through expected correlations with related constructs. When the instrument will be used to compare groups or predict outcomes, measurement invariance testing becomes essential \citep{AERA2014Standards}.

\hypertarget{construct-irrelevant-variance-in-educational-assessment}{%
\subsubsection{2.3 Construct-Irrelevant Variance in Educational Assessment}\label{construct-irrelevant-variance-in-educational-assessment}}

The unified validity framework identifies construct-irrelevant variance as a principal threat to score interpretation \citep{AERA2014Standards,Messick1995Validity}. CIV occurs when systematic variance in test scores is attributable to factors extraneous to the target construct. When assessment content intersects with respondents' attitudes or social identities, CIV can arise through attitude-consistent responding: students whose beliefs align with test content may find certain response options more natural or appealing, inflating their scores without corresponding gains in the assessed competence \citep{Messick1995Validity}.

For the HPT instrument specifically, the concern is that the Weimar Republic scenario creates a \emph{congruence pathway}---a mechanism whereby attitudinal similarity to item content lowers the cognitive effort required to endorse a response, inflating scores independently of competence. Consider a contextualization item such as ``His decision was understandable given the political and economic turmoil of the Weimar Republic.'' A student who endorses this statement because she has reconstructed the historical context produces a valid HPT response. But a student who endorses it because the sentiment resonates with her own political sympathies---not because she has engaged in historical reasoning---produces an inflated score. This would represent construct contamination in \citeauthor{Messick1995Validity}'s (\citeyear{Messick1995Validity}) terms---a systematic source of score variance unrelated to HPT competence. Evidence of internal structure, including measurement invariance across relevant subgroups, constitutes essential validity evidence when scores are used to compare groups or when the construct definition requires demonstrating that extraneous factors do not influence scores \citep{AERA2014Standards}.

An alternative theoretical prediction deserves mention. Research on authoritarian cognition suggests that high authoritarianism is associated with reduced integrative complexity---the cognitive capacity to hold and reconcile multiple conflicting perspectives simultaneously \citep{Duckitt2001,Suedfeld1977}. Because HPT requires precisely this kind of integrative capacity, authoritarian students might perform \emph{worse}, not better. The direction of any bias, if present, was genuinely uncertain. Additionally, dual-process theories of cognition \citep{Evans2013} suggest that structured assessment formats may channel deliberative processing, overriding the automatic attitude-retrieval processes that would produce attitude-consistent responding. If so, the assessment format itself may buffer against attitudinal interference.

\hypertarget{dif-and-measurement-invariance-as-fairness-evidence}{%
\subsubsection{2.4 DIF and Measurement Invariance as Fairness Evidence}\label{dif-and-measurement-invariance-as-fairness-evidence}}

Differential item functioning and measurement invariance testing provide complementary approaches to evaluating measurement fairness \citep{Chen2007,Cheung2002,Finch2005}. DIF analysis examines whether individual items function differently across groups after controlling for the latent trait, identifying specific items that may be biased. Measurement invariance testing through multi-group CFA examines whether the overall measurement model holds across groups at three progressively restrictive levels: \emph{configural invariance} (the same factors exist in both groups), \emph{metric invariance} (items relate equally strongly to the factors in both groups), and \emph{scalar invariance} (the scale origin is identical, so mean scores can be directly compared across groups).

For ordinal items, the weighted least squares mean and variance
adjusted (WLSMV) estimator---a robust estimation method designed for
categorical data---is recommended for CFA \citep{Flora2004}.
Because ordinal items have discrete response categories rather than
continuous scales, threshold-based invariance testing (which models the
boundaries between response categories) replaces intercept-based
approaches. The conventional criteria for evaluating invariance steps use changes in fit indices rather than chi-square difference tests, which are sensitive to sample size: \(\Delta\)CFI \(\leq\) .01 and \(\Delta\)RMSEA \(\leq\) .015 indicate that the more constrained model does not fit meaningfully worse \citep{Chen2007,Cheung2002}.

In the context of HPT measurement, establishing scalar invariance across ideology groups would mean that the factor structure, loadings, and item thresholds are equivalent for students with different political attitudes---a necessary condition for concluding that observed score differences (or the lack thereof) reflect true construct differences rather than measurement artifacts.

\hypertarget{the-present-study}{%
\subsubsection{2.5 The Present Study}\label{the-present-study}}

This study addresses two interrelated measurement questions:

\begin{enumerate}
\def\labelenumi{\arabic{enumi}.}
\item
  What psychometric evidence does the Czech adaptation of the Hartmann and Hasselhorn HPT instrument provide, including evidence concerning a correlated three-factor representation, score reliability, and convergence with historical knowledge?
\item
  What evidence do DIF, measurement-invariance, and regression analyses provide about whether scores are affected by students' political attitudes?
\end{enumerate}

The preregistration specified six hypotheses: positive ideology
associations with CONT and total HPT, a negative association with raw
POP, positive DIF on CONT items, stronger ideology--HPT associations at
lower knowledge, and exploratory correlation contrasts. It did not
predict invariance. The present measurement-focused organization and
several analyses---including HPT\_CTX6 as the primary composite,
all-item GRM DIF, the bifactor model, NS-only invariance, and
equivalence tests---were adopted after registration. Not all registered
analyses were implemented. Table S5 maps the registered hypotheses,
omissions, and post-registration additions; claims from the latter are
treated as exploratory or descriptive rather than preregistered tests.

\hypertarget{method}{%
\subsection{3. Method}\label{method}}

\hypertarget{adaptation-process}{%
\subsubsection{3.1 Adaptation Process}\label{adaptation-process}}

The Czech HPT instrument was developed through a systematic translation and adaptation process following ITC guidelines \citep{ITC2017}. Two bilingual Czech-English experts conducted independent forward translations from the German instrument reported in \citet{Hartmann2008}; wording was cross-checked against the published English validation article \citep{Hartmann_Hasselhorn2008Historical}. An independent backward translation was produced by a third bilingual expert unfamiliar with the original. Five Czech history educators at the secondary level reviewed the adapted items for conceptual equivalence, age-appropriateness, and curricular alignment. No items required substantive modification; minor wording adjustments were made to two items for naturalness in Czech. Reviewers judged the scenario content accessible to Czech secondary students given the curriculum coverage described in Section 2.2. This was expert-review evidence, not direct evidence from student interviews.

The adaptation preserved the four-point response format (1 = \emph{strongly disagree} to 4 = \emph{strongly agree}), the three item categories (three items per category), and reverse-scoring of POP items from the German source. All deviations from the original were documented transparently in the OSF repository (https://doi.org/10.17605/OSF.IO/YNG37). The review panel cross-checked the scenario against the Czech history curriculum, but no direct student response-process study was conducted.

No cognitive pilot study (e.g., think-aloud interviews) was conducted prior to main data collection, which constitutes a limitation of the adaptation process (see Section 5.6).

\hypertarget{participants}{%
\subsubsection{3.2 Participants}\label{participants}}

The sample comprised 293 students from 20 classrooms across 10
schools in the Czech Republic, spanning lower-secondary
(International Standard Classification of Education [ISCED] level 2,
approximately ages 11--15; \emph{n} = 87, 29.7\%) and upper-secondary
(ISCED 3, approximately ages 15--19; \emph{n} = 206, 70.3\%) levels. The inclusion criterion was completion of the curricular unit on the interwar period. Schools were drawn from five Czech regions (Praha: 42.7\%; Královéhradecký: 34.8\%; Zlínský: 11.3\%; Jihomoravský: 8.2\%; Ústecký: 3.1\%) and three funding categories (public: 90.8\%; private: 6.5\%; church-affiliated: 2.7\%). Gender composition was 192 male (65.5\%), 88 female (30.0\%), and 13 identifying as another gender (4.4\%). Students' most recent history grade averaged 1.71 (\emph{SD} = 0.92; Czech scale: 1 = excellent, 5 = failing).

Schools were recruited through convenience sampling; the sample over-represents upper-secondary students relative to the national age cohort composition and should not be treated as a probability sample of Czech adolescents. Sample size was determined by an a priori power analysis specified in the preregistered protocol (OSF: https://doi.org/10.17605/OSF.IO/YNG37), targeting 80\% power to detect standardized ideology effects of \(\beta\) \(\geq\) .15 in a two-level multilevel regression; the achieved \emph{N} = 285--293 met that regression target. The calculation did not establish power for the post-registration DIF or MG-CFA analyses. Within each school, all participating students present on the day of administration completed the instrument battery during a regular class period (approximately 30--40 minutes). Teachers were present during administration but did not intervene in students' responses.

\hypertarget{instruments}{%
\subsubsection{3.3 Instruments}\label{instruments}}

\hypertarget{historical-perspective-taking-hpt}{%
\paragraph{3.3.1 Historical Perspective Taking (HPT)}\label{historical-perspective-taking-hpt}}

Students completed the Czech adaptation of the scenario-based HPT instrument from the German dissertation \citep{Hartmann2008}, with the published English account used as supporting validation evidence \citep{Hartmann_Hasselhorn2008Historical}. They read a narrative in which Hannes discusses Germany's crisis and the 1930 election, then judged nine explanations of his situation, including whether he could vote for an anti-democratic party such as the NSDAP, on a 4-point Likert scale. Three items each tapped present-oriented perspective (POP1--POP3), role of the historical agent (ROA1--ROA3), and historical contextualization (CONT1--CONT3). POP items capture present-centered judgments of the historical figure; these are reverse-scored (5 $-$ POP, yielding a reversed range of 1--4) so that higher scores indicate less present-oriented reasoning. ROA items explain action through familiar or stereotyped social roles. CONT items require students to situate the historical figure's deliberation within the specific political, economic, and social conditions of the Weimar Republic. The primary HPT\_CTX6 composite was the mean of reversed POP and CONT scores and therefore represents the POP--CONT contextualization dimension. HPT\_TOT9 (adding the theoretically less certain ROA category) served only as a secondary sensitivity outcome.

\hypertarget{right-authoritarian-ideology-fr-lf-mini}{%
\paragraph{3.3.2 Right-Authoritarian Ideology (FR-LF mini)}\label{right-authoritarian-ideology-fr-lf-mini}}

A six-item short form of the Fragebogen Rechtsextremer Einstellungen---Leipziger Form (Questionnaire on Right-Wing Extremist Attitudes---Leipzig Form; \citealp{Decker_et_al2013Rechtsextremismus}), comprising acceptance of right-wing dictatorship (RD; 3 items) and National Socialist sympathy/relativization (NS; 3 items) facets, rated on a 5-point scale. The NS facet is theoretically the most direct operationalization of potential content congruence with the Weimar scenario.

\hypertarget{authoritarianism-ksa-3}{%
\paragraph{3.3.3 Authoritarianism (KSA-3)}\label{authoritarianism-ksa-3}}

The nine-item Kurzskala Autoritarismus (Short Scale of
Authoritarianism; \citealp{Beierlein_et_al2014Die_Kurzskala}), measuring authoritarian
aggression, submission, and conventionalism on a 5-point scale. This instrument was used for a supplementary measurement invariance check using an alternative ideology operationalization.

\hypertarget{historical-knowledge-kn}{%
\paragraph{3.3.4 Historical Knowledge (KN)}\label{historical-knowledge-kn}}

Six multiple-choice items covering content relevant to the HPT scenario, scored dichotomously (sum score 0--6). This measure served as a convergent validity criterion, given the theoretical expectation that HPT is a knowledge-dependent competence.

\hypertarget{social-desirability-sdr-5}{%
\paragraph{3.3.5 Social Desirability (SDR-5)}\label{social-desirability-sdr-5}}

The five-item short form \citep{Hays_et_al1989A_Five}, rated on a 5-point scale with items SDR2--SDR4 reverse-coded. Included to evaluate whether HPT scores reflect impression management rather than genuine reasoning. The source studies for the FR-LF included German respondents aged 14 years and older, whereas KSA-3 was developed for German-speaking adults aged 18 years and older and SDR-5 with adult medical outpatients and older adults. None therefore constitutes direct Czech adolescent validation; suitability rests on expert review and the sample-specific evidence reported here and in Table S4.

\hypertarget{analytic-approach}{%
\subsubsection{3.4 Analytic Approach}\label{analytic-approach}}

All analyses were conducted in R (version 4.4.2) using lavaan for CFA, mirt for IRT-based DIF analysis, psych for reliability estimation, and lme4/performance for intraclass correlations.

\textbf{Factor structure.} CFA compared one-factor (general HPT), two-factor (POP+CONT vs.~ROA), and three-factor (POP, ROA, CONT) models using the WLSMV estimator, appropriate for ordinal item responses \citep{Flora2004}. CFA fit indices summarize how well a hypothesized factor structure reproduces the observed correlations among items; multiple indices are reported because each captures a different aspect of misfit. Model evaluation used CFI (\(\geq\) .95 acceptable, \(\geq\) .97
good), the Tucker--Lewis Index (TLI; \(\geq\) .95), the Root Mean
Square Error of Approximation (RMSEA; \(\leq\) .08 acceptable,
\(\leq\) .05 good), and the Standardized Root Mean Square Residual
(SRMR; \(\leq\) .08).

\textbf{Reliability.} Internal consistency was assessed using raw and
polychoric Cronbach's \(\alpha\), McDonald's \(\omega\), and KR-20 for
the dichotomous knowledge items. Reliability of the combined ideology
score was estimated with the Mosier formula. Full methods and estimates
are reported in Table S3.

\textbf{Intraclass correlations.} We computed intraclass correlation
coefficients (ICCs) at the school and classroom levels to characterize
the multilevel structure of the data and to evaluate whether
clustering warranted multilevel approaches.

\textbf{Clustered data handling.} The substantive regressions used
random intercepts for the nested data. The DIF, MG-CFA, and zero-order
correlation analyses treated observations as independent. Composite
score ICCs do not establish that item thresholds, loadings, or their
standard errors are unaffected by clustering; therefore, uncertainty
for these analyses may be understated. The zero-order correlation
intervals are reported as descriptive, unadjusted intervals rather than
cluster-robust inference.

\textbf{Missing data and score construction.} Composite scores required
at least two answered items per three-item HPT
subscale, four of six FR-LF items, seven of nine KSA-3 items, and four
of five SDR-5 items; available responses were averaged. Missing
knowledge responses were scored as incorrect, following the original
analysis scripts. Reliability estimates used listwise-complete item
sets, while correlations used pairwise-complete composite scores. These
rules explain the analysis-specific sample sizes reported in the tables
and supplement.

\textbf{Differential item functioning.} Low and High ideology groups were formed post-registration using tertile splits on a composite ideology score (mean of \emph{z}-scored FR-LF and KSA-3), with the middle tertile excluded (\emph{n} = 96 per group). A constrained unidimensional multi-group graded response model (GRM item
discrimination and category-intercept parameters for each item by group are
reported in Table S1 in the OSF supplementary materials) using a
sequential item-release design tested each item for omnibus DIF by
jointly freeing its discrimination and category intercepts, with
Bonferroni correction (\(\alpha\) = .01). The
baseline model constrains all nine items to equality across groups;
each test then frees the discrimination and category intercepts for
the focal item, yielding a likelihood-ratio comparison against the
fully constrained baseline (an all-others-as-anchor strategy). This
approach assumes that the remaining eight items are DIF-free when
testing each focal item; iterative purification was not applied, which
means that if multiple anchor items exhibited small DIF, focal-item
detection could be affected.
Because the CFA favoured correlated dimensions, the unidimensional GRM
is a pragmatic sensitivity model rather than a fully aligned
measurement model; multidimensionality or local dependence may affect
the item tests.

\textbf{Sensitivity of DIF detection.} With approximately 96 respondents per ideology group and nine Bonferroni-adjusted tests, statistical power was limited. No design-specific simulation was conducted, so the minimum detectable DIF magnitude is unknown. The extreme-group design increases ideological contrast at the cost of excluding the middle tertile and reducing group size.

\textbf{Multi-group CFA.} Stepwise MG-CFA tested configural, metric, and scalar invariance across ideology groups using the three-factor model with WLSMV estimation. Invariance was evaluated using \(\Delta\)CFI \(\leq\) .01 and \(\Delta\)RMSEA \(\leq\) .015 criteria \citep{Chen2007,Cheung2002}.

\textbf{Convergent validity.} Descriptive Pearson correlations between HPT composites and historical knowledge provided convergent validity evidence; Fisher-transformed intervals did not adjust for clustering.

\textbf{Supplementary invariance check.} To verify that invariance findings were not an artifact of the composite ideology definition, MG-CFA was also conducted using NS-only grouping (the facet most directly content-congruent with the Weimar scenario).

\hypertarget{results}{%
\subsection{4. Results}\label{results}}

\hypertarget{factor-structure}{%
\subsubsection{4.1 Factor Structure}\label{factor-structure}}

CFA supported the three-factor model (POP, ROA, CONT) as best-fitting, outperforming the two-factor and one-factor alternatives (Table 1). The three-factor model showed excellent global fit: CFI = .999, TLI = .998, RMSEA = .011 (90\% CI [.000, .051]), SRMR = .048. The two-factor model combining POP and CONT items on a single factor showed acceptable but inferior fit (CFI = .956, RMSEA = .057), while the one-factor model fit poorly (CFI = .926, RMSEA = .073).

\textbf{Table 1.} \emph{CFA Fit Indices for HPT Factor Models (WLSMV Estimator)}

\begin{longtable}[]{@{}lrrrrrll@{}}
\toprule
Model & \(\chi^2\) & \emph{df} & CFI & TLI & RMSEA & 90\% CI & SRMR \\
\midrule
\endhead
\bottomrule
\endlastfoot
1-Factor & 66.68 & 27 & .926 & .902 & .073 & [.051, .095] & .075 \\
2-Factor (POP+CONT, ROA) & 49.48 & 26 & .956 & .940 & .057 & [.032, .081] & .066 \\
3-Factor (POP, ROA, CONT) & 24.78 & 24 & .999 & .998 & .011 & [.000, .051] & .048 \\
\end{longtable}

\emph{Note.} WLSMV = weighted least squares mean and variance adjusted. All items treated as ordered categorical. The 2-factor model groups POP and CONT items on a single factor with ROA on a separate factor. Chi-square values are mean-and-variance adjusted (WLSMV). Nested model comparisons require the DIFFTEST procedure rather than standard chi-square difference tests. All \emph{p} $<$ .001 except the three-factor model (\emph{p} = .418).

The near-perfect global fit of the three-factor model reflects the
small item set (nine items in three just-identified three-indicator
factors---each subscale has exactly as many free parameters as observed
statistics, leaving no room for misfit within subscales) and should be
interpreted alongside local fit diagnostics: the largest modification
index was 7.35 (POP1--ROA3 residual covariance; below the threshold
of 10), and standardized residual correlations were small (\(|r|\)
\(\leq\) .13). With only three indicators per factor, there is limited
opportunity to reveal within-factor misfit; the near-perfect CFI is therefore a
structural consequence of the low-\emph{df} specification rather than
evidence of an unusually well-fitting model. The Czech data are
compatible with the correlated three-factor specification, but this is
not a test of cross-national equivalence.

Inter-factor correlations from the three-factor CFA were: POP--ROA \(\Phi\) = .377, POP--CONT \(\Phi\) = .509, ROA--CONT \(\Phi\) = .631. These moderate-to-strong correlations and the better fit of the three-factor model provide internal-structure evidence that the item sets are distinguishable in this sample. They do not by themselves establish discriminant validity. Stronger evidence would require showing that the subscales relate differentially to external criteria in theoretically expected ways.

A supplementary bifactor model (one general factor plus three specific
factors, in which each item loads on both the general factor and its
subscale-specific factor) fit well (\(\chi^2\) = 14.70, \emph{df} =
18, \emph{p} = .682) but did not meaningfully improve over the
correlated three-factor model (\(\Delta\)CFI = .001). The low
\emph{df} (18) results from the additional cross-loadings consuming
degrees of freedom, and the comparison with the correlated three-factor
model (\emph{df} = 24) should be treated cautiously given the
different parameterisations. CONT items loaded primarily on the
general factor (.47--.63), while POP and ROA retained meaningful
specific factor variance, supporting subscale-level interpretation.

\hypertarget{standardized-factor-loadings}{%
\subsubsection{4.2 Standardized Factor Loadings}\label{standardized-factor-loadings}}

Table 2 presents the standardized loadings from the three-factor CFA. CONT items loaded most strongly (.61--.72), followed by POP items (.30--.74) and ROA items (.30--.67). POP2 showed the weakest loading (.30), at the conventional minimum threshold for item retention; it was retained to preserve fidelity to the original instrument structure.

\textbf{Table 2.} \emph{Summary of Standardized Factor Loadings}

\begin{longtable}[]{@{}lll@{}}
\toprule
Subscale & Loading range & Weakest item \\
\midrule
\endhead
\bottomrule
\endlastfoot
Present-oriented perspective & .30--.74 & POP2 (.30) \\
Role of historical agent & .30--.67 & ROA2 (.30) \\
Historical contextualization & .61--.72 & CONT2 (.61) \\
\end{longtable}

\emph{Note.} POP items are reverse-scored (higher = more contextualized reasoning). All loadings were significant at \emph{p} $<$ .001. Full item estimates, standard errors, and confidence intervals are reported in Table S2.

\hypertarget{reliability-and-internal-consistency}{%
\subsubsection{4.3 Reliability and Internal Consistency}\label{reliability-and-internal-consistency}}

Table 3 summarizes reliability for the primary study scores. Across HPT subscales, raw Cronbach's \(\alpha\) ranged from .46 to .64 and McDonald's \(\omega\) from .54 to .70. Full estimates for all HPT and predictor scores are in Table S3.

\textbf{Table 3.} \emph{Reliability Summary for Primary Study Scores}

\begin{longtable}[]{@{}lccc@{}}
\toprule
Score & Items/components & Reliability & Method \\
\midrule
\endhead
\bottomrule
\endlastfoot
HPT composite & 6 & .66 & \(\omega\) total \\
HPT total & 9 & .70 & \(\omega\) total \\
Combined ideology & 2 scales & .80 & Mosier \\
\end{longtable}

\emph{Note.} The Mosier estimate has a 95\% bootstrap CI [.74, .84]. Full subscale and predictor alpha, polychoric-alpha, omega, and interval estimates are reported in Table S3.

Reliability was modest and constrains individual interpretation. FR-LF
had \(\alpha\) = .67, KSA-3 had \(\alpha\) = .72, and the combined
ideology score had estimated reliability .80, 95\% CI [.74, .84].
Knowledge had KR-20 = .55, and SDR-5 was weak (\(\alpha\) = .30).

\hypertarget{descriptive-statistics}{%
\subsubsection{4.4 Descriptive Statistics}\label{descriptive-statistics}}

Table 4 presents descriptive statistics. HPT scores followed the expected mean-score ordering (POP\_rev $>$ ROA $>$ CONT), with contextualization producing the lowest scores (\emph{M} = 2.71); this ordering may reflect item endorsement or difficulty and is not developmental evidence. Ideology scores were below the scale midpoint (FR-LF: \emph{M} = 2.48; KSA: \emph{M} = 2.86). The FR-LF median was 2.33 (IQR [2.00, 3.00]) and 51 of 284 students (18.0\%) scored above its midpoint. The KSA-3 median was 2.89 (IQR [2.53, 3.22]) and 106 of 283 students (37.5\%) scored above its midpoint. On the mean of the two raw scales, 71 of 283 students (25.1\%) scored above 3. Historical knowledge had a median of 3 (IQR [2, 4]), with 4.1\% at the floor and 8.9\% at the ceiling. No defensible student- or school-level socioeconomic measure was available.

\textbf{Table 4.} \emph{Descriptive Statistics for Study Variables}

\begin{longtable}[]{@{}llllll@{}}
\toprule
Variable & \emph{n} & \emph{M} & \emph{SD} & Min & Max \\
\midrule
\endhead
\bottomrule
\endlastfoot
HPT POP (reversed) & 287 & 2.97 & 0.64 & 1.33 & 4.00 \\
HPT ROA & 287 & 2.80 & 0.68 & 1.00 & 4.00 \\
HPT CONT & 287 & 2.71 & 0.74 & 1.00 & 4.00 \\
HPT CTX6 (primary) & 287 & 2.84 & 0.55 & 1.33 & 4.00 \\
HPT TOT9 & 287 & 2.83 & 0.49 & 1.67 & 3.89 \\
KN Total (0--6) & 293 & 3.04 & 1.62 & 0.00 & 6.00 \\
FR-LF RD & 285 & 2.54 & 0.88 & 1.00 & 5.00 \\
FR-LF NS & 286 & 2.43 & 0.89 & 1.00 & 5.00 \\
FR-LF Total & 284 & 2.48 & 0.75 & 1.00 & 5.00 \\
KSA-3 Total & 283 & 2.86 & 0.62 & 1.00 & 4.60 \\
SDR-5 Total & 283 & 3.02 & 0.63 & 1.00 & 4.60 \\
\end{longtable}

\emph{Note.} HPT items scored on a 1--4 scale; POP items reverse-coded. FR-LF and KSA-3 scored on a 1--5 scale. KN scored as sum of correct responses (0--6).

\hypertarget{intraclass-correlations}{%
\subsubsection{4.5 Intraclass Correlations}\label{intraclass-correlations}}

Intraclass correlations indicated meaningful school-level clustering for most HPT outcomes (HPT\_CTX6: ICC\_school = .041; HPT\_TOT9: ICC\_school = .085; ROA: ICC\_school = .102), consistent with between-school differences in instructional emphasis and student composition. Classroom-level variance was near zero for the primary composite but non-negligible for CONT (ICC\_class = .033) and ROA (ICC\_class = .104). These ICCs justify multilevel modelling for substantive analyses.
The school-level variance is consistent with between-school
differences in history instruction, curricular emphasis, and student
composition, although the present design cannot identify which of
these mechanisms produced the variation.

\hypertarget{dif-analysis}{%
\subsubsection{4.6 DIF Analysis}\label{dif-analysis}}

Table 5 presents omnibus DIF results for all nine HPT items. Using a constrained multi-group graded response model with likelihood-ratio tests and Bonferroni-adjusted significance (\(\alpha\) = .01), no items were flagged. All adjusted \emph{p}-values exceeded .01. No DIF was detected between the low- and high-ideology groups with this sample and procedure; the magnitude that the design could reliably detect is unknown.

\textbf{Table 5.} \emph{DIF Results per HPT Item (Likelihood-Ratio Tests from Multi-Group Graded Response Model)}

\begin{longtable}[]{@{}llllll@{}}
\toprule
Item & \(\chi^2\) & \emph{df} & \emph{p} (raw) & \emph{p} (Bonferroni) & Flagged \\
\midrule
\endhead
\bottomrule
\endlastfoot
POP1 & 11.18 & 4 & .025 & .222 & No \\
POP2 & 9.42 & 4 & .051 & .462 & No \\
POP3 & 9.50 & 4 & .050 & .448 & No \\
ROA1 & 9.88 & 4 & .043 & .383 & No \\
ROA2 & 2.90 & 4 & .575 & 1.000 & No \\
ROA3 & 8.03 & 4 & .091 & .815 & No \\
CONT1 & 1.65 & 4 & .800 & 1.000 & No \\
CONT2 & 3.45 & 4 & .486 & 1.000 & No \\
CONT3 & 9.43 & 4 & .051 & .461 & No \\
\end{longtable}

\emph{Note.} Likelihood-ratio chi-square statistics from constrained
multi-group graded response models. Each test compares a fully
constrained baseline (all items equal across groups) to a model
freeing discrimination and category-intercept parameters for the tested item.
Low (\emph{n} = 96) and High (\emph{n} = 96) ideology groups formed
by tertile split on composite ideology score (middle tertile
excluded). Bonferroni-adjusted \(\alpha\) = .01 (9 tests). 0 of 9
items flagged. POP1 showed the largest DIF statistic (\(\chi^2\) =
11.18, \emph{p}\_raw = .025), well above the Bonferroni threshold
after correction (\emph{p}\_adj = .222). GRM item discrimination and
category-intercept parameters are reported in Table S1 in the OSF
supplementary materials.

Because every other item served as an anchor, undetected DIF among
those anchors could affect each test. The result should therefore be
interpreted as ``no DIF detected with this sample and procedure,'' not
``no DIF present.''

\hypertarget{measurement-invariance}{%
\subsubsection{4.7 Measurement Invariance}\label{measurement-invariance}}

Table 6 presents the stepwise MG-CFA results. The configural model produced CFI = .978, RMSEA = .043, and SRMR = .074. Adding equality constraints did not worsen CFI or RMSEA: metric \(\Delta\)CFI = .000 and \(\Delta\)RMSEA = $-$.002; scalar \(\Delta\)CFI = +.002 and \(\Delta\)RMSEA = $-$.006. These changes fall within common fit-index criteria \citep{Chen2007,Cheung2002}. However, the configural and metric solutions were inadmissible, so the sequential invariance hierarchy was not established and the fit-index changes are reported descriptively only.

\textbf{Table 6.} \emph{Multi-Group CFA Fit Indices Across Invariance Levels (WLSMV)}

\begin{longtable}[]{@{}llllll@{}}
\toprule
Model & CFI & RMSEA & SRMR & \(\Delta\)CFI & \(\Delta\)RMSEA \\
\midrule
\endhead
\bottomrule
\endlastfoot
Configural & .978 & .043 & .074 & --- & --- \\
Metric & .978 & .041 & .081 & .000 & $-$.002 \\
Scalar & .980 & .035 & .076 & +.002 & $-$.006 \\
\end{longtable}

\emph{Note.} \(\Delta\)CFI and \(\Delta\)RMSEA computed relative to the preceding model. Invariance criteria: $|\Delta\text{CFI}|$ \(\leq\) .01, $|\Delta\text{RMSEA}|$ \(\leq\) .015 \citep{Chen2007}. Low and High ideology groups formed by tertile split on composite ideology (\emph{n} \(\approx\) 96 per group). Items treated as ordered categorical.

The SRMR increased slightly from configural (.074) to metric (.081),
marginally exceeding the conventional \(\leq\) .08 threshold. This
threshold was developed for maximum-likelihood estimation with
continuous indicators and may not apply directly to WLSMV-estimated
models with polychoric residuals; the marginal exceedance should
therefore be interpreted cautiously. The SRMR returned to .076 at the
scalar step. This pattern can signal partial loading non-invariance.
The largest modification index at the metric step was 11.28 (a
suggested cross-loading of POP2 on the ROA factor in the Low ideology
group), marginally above the conventional threshold of 10. This
suggests a minor loading difference for POP2 that is consistent with
its weak loading (.30) in the overall CFA (Table 2) rather than
evidence of systematic group-specific misspecification.

Negative estimated residual variances (Heywood cases) occurred for ROA
items in the Low ideology group in the configural and metric models.
Such estimates indicate a non-positive-definite covariance matrix,
which can invalidate fit statistics and render parameter estimates
uninterpretable. While Heywood cases can arise with small per-group
samples and WLSMV estimation of ordinal items---particularly with
three-item factors---they signal potential model inadmissibility rather
than a benign estimation artefact. The per-group sample sizes
(\emph{n} \(\approx\) 82--96) were small for this model, particularly
with three-item factors. Consequently, no configural, metric, or scalar
invariance level can be
firmly established from this sequence. The fit-index pattern is
compatible with equality constraints, but the inadmissible baseline
prevents it from supporting scalar equivalence.

\hypertarget{supplementary-invariance-ns-only-grouping}{%
\subsubsection{4.8 Supplementary Invariance: NS-Only Grouping}\label{supplementary-invariance-ns-only-grouping}}

To verify that invariance findings were not an artefact of the
composite ideology grouping variable, we repeated the MG-CFA using
groups defined by the NS facet alone (National Socialist
sympathy/relativization), which represents the most direct theoretical
operationalization of content congruence with the Weimar scenario. With
the NS tertile split, group sizes were \emph{n} = 126 (Low) and
\emph{n} = 125 (High), with 42 middle-group cases excluded. Fit did not
deteriorate when equality constraints were added (configural CFI = .960,
RMSEA = .050, SRMR = .070;
metric \(\Delta\)CFI = +.002, \(\Delta\)SRMR = +.006; scalar
\(\Delta\)CFI = +.014, \(\Delta\)SRMR = $-$.003). Because CFI
increased and RMSEA and SRMR decreased at the scalar step, the positive
\(\Delta\)CFI does not indicate worse fit and should not be compared
with a deterioration threshold by absolute magnitude. This
post-registration analysis provides descriptive robustness evidence,
not a second confirmatory test. Full fit statistics for this
supplementary analysis are available in the OSF repository.

\hypertarget{convergent-validity}{%
\subsubsection{4.9 Convergent Validity}\label{convergent-validity}}

Scenario-relevant historical knowledge correlated positively with all
HPT outcomes: \emph{r} = .34 (HPT\_CTX6), \emph{r} = .22 (CONT),
\emph{r} = .32 (POP\_rev). These moderate correlations indicate that
students with stronger knowledge of the relevant historical period
scored higher on HPT. The scores were not empirically interchangeable,
but their unshared variance cannot be attributed specifically to
reasoning because it also includes measurement error and other
unmeasured influences. A limitation of this evidence is that
the KN items were written to cover content directly relevant to the
HPT scenario, so the correlation may partly reflect shared
specific-content activation rather than the broader
``knowledge-dependent competence'' relationship posited theoretically.
A general historical knowledge measure would provide a stronger
convergent validity test.

History grade also correlated with HPT\_CTX6 (\emph{r} = $-$.153, \emph{p} = .010; negative sign reflects the Czech scale where 1 = excellent). By contrast, ideology measures showed negligible correlations with HPT (FR-LF--HPT\_CTX6: \emph{r} = $-$.05), and social desirability was unrelated (\emph{r} = $-$.02 to .01). This pattern---significant knowledge correlations, null ideology and social desirability correlations---is consistent with an instrument that measures disciplinary reasoning rather than attitudes or impression management.

\textbf{Table 7.} \emph{Zero-Order Pearson Correlations Among Key Variables}

\begin{scriptsize}
\begin{longtable}[]{@{}p{2.7cm}rrrrrrr@{}}
\toprule
 & HPT & Context & Pres. (rev.) & Knowledge & FR-LF & KSA-3 & SDR-5 \\
\midrule
\endhead
\bottomrule
\endlastfoot
HPT composite & --- & [.78, .86] & [.71, .81] & [.23, .43] & [$-$.16, .07] & [$-$.15, .09] & [$-$.14, .10] \\
Contextualization & .82*** & --- & [.14, .36] & [.11, .33] & [$-$.10, .13] & [$-$.05, .19] & [$-$.11, .13] \\
Present-oriented perspective (rev.) & .76*** & .26*** & --- & [.20, .41] & [$-$.21, .02] & [$-$.25, $-$.02] & [$-$.16, .07] \\
Knowledge & .34*** & .23*** & .31*** & --- & [$-$.22, .01] & [$-$.07, .16] & [$-$.09, .15] \\
FR-LF & $-$.05 & .02 & $-$.10 & $-$.10 & --- & [.43, .60] & [$-$.30, $-$.08] \\
KSA-3 & $-$.03 & .07 & $-$.13* & .05 & .52*** & --- & [$-$.26, $-$.03] \\
SDR-5 & $-$.02 & .01 & $-$.05 & .03 & $-$.19** & $-$.15* & --- \\
\end{longtable}
\end{scriptsize}

\emph{Note.} Correlations are below the diagonal and unadjusted descriptive 95\% confidence intervals are above it. Pairwise complete observations (\emph{n} = 282--287; exact values in Table S3b). Present-oriented perspective is reverse-scored. Knowledge = historical knowledge (0--6 sum); FR-LF = right-authoritarian ideology (1--5); KSA-3 = authoritarianism (1--5); SDR-5 = social desirability (1--5). Intervals and \emph{p}-values do not adjust for school or classroom clustering. *** \emph{p} $<$ .001. ** \emph{p} $<$ .01. * \emph{p} $<$ .05.

Figure~\ref{fig:relationships} visualizes the near-zero ideology
association and positive knowledge association with HPT\_CTX6.

\begin{figure}[htbp]
\centering
\includegraphics[width=\textwidth]{revision_outputs/revision_relationships.pdf}
\caption{Observed relationships of HPT\_CTX6 with the combined ideology score---the mean of standardized FR-LF and KSA-3 scores---(Panel A) and historical knowledge (Panel B). Points are jittered to display overlap; lines and shaded regions show unadjusted linear fits and 95\% confidence intervals.}
\label{fig:relationships}
\end{figure}

\hypertarget{substantive-context-ideology-hpt-relationships}{%
\subsubsection{4.10 Substantive Context: Ideology-HPT Relationships}\label{substantive-context-ideology-hpt-relationships}}

Ideology was unrelated to HPT scores across all model
specifications: FR-LF \(\beta\) = 0.012, \emph{SE} = 0.067, 95\%
CI [$-$0.120, 0.143], \emph{p} = .863 for the primary outcome
(HPT\_CTX6). Historical knowledge was the only consistent
predictor (\(\beta\) = 0.21--0.34, \emph{p} $<$ .001).
As an exploratory sensitivity analysis, equivalence tests used the Two
One-Sided Tests (TOST) procedure, which reverses the logic of
null-hypothesis testing by asking whether an effect is small enough to
be declared practically negligible. The Smallest Effect Size Of
Interest (SESOI) was set post hoc at \(\pm\) 0.20 standardized units,
corresponding to \citeauthor{Cohen1988}'s (\citeyear{Cohen1988}) threshold for a ``small'' effect.
This SESOI is wider than the preregistered power target of \(\beta\)
\(\geq\) .15, meaning that effects between .15 and .20 cannot be
confidently excluded; the TOST results should therefore be treated as
exploratory rather than confirmatory. With the \(\pm\) 0.20 boundary,
the TOST results indicated that ideology effects at or above \(\beta\)
= .20 are statistically ruled out (all TOST \emph{p}s $<$ .003; 90\%
CI for the primary outcome [$-$0.10, 0.12], entirely within the
\(\pm\) 0.20 SESOI). The ideology coefficient remained non-significant
across all sensitivity checks, including alternative HPT scoring
composites, different ideology operationalizations, exclusion of
high social desirability responders, random-slope models, and
school fixed-effects specifications (all \emph{p}s $>$
.43). These substantive findings complement the measurement
evidence: no DIF or composite-level ideology
association was detected within the observed ideological range.

\hypertarget{discussion}{%
\subsection{5. Discussion}\label{discussion}}

\hypertarget{summary-of-findings}{%
\subsubsection{5.1 Summary of Findings}\label{summary-of-findings}}

The results provide qualified fairness evidence. The expected factor
structure was supported, no DIF was detected with this procedure, and
ideology did not predict HPT scores; scalar-invariance evidence was
qualified by inadmissible smaller-group solutions. The Czech adaptation showed modest subscale reliability, a positive relationship with scenario-relevant historical knowledge, and sensitivity to school-level variation. Political ideology, authoritarianism, and social desirability were unrelated to HPT, although weak reliability constrains the latter null associations.

\hypertarget{factor-structure-evidence}{%
\subsubsection{5.2 Factor-Structure Evidence}\label{factor-structure-evidence}}

The correlated three-factor model (POP, ROA, CONT) fit the Czech data better than the tested alternatives. This does not demonstrate cross-cultural equivalence, and the specification differs from the two-dimensional PCA pattern in the original German study. It shows only that the Czech data are compatible with separating the three item sets, under a near-saturated model with limited opportunity to reveal misfit. CONT items loaded most consistently (.61--.72), while POP2 and ROA2 showed weaker loadings (.30); these were retained for comparability but are candidates for revision in future adaptations.

\hypertarget{addressing-modest-reliability}{%
\subsubsection{5.3 Addressing Modest Reliability}\label{addressing-modest-reliability}}

The subscale reliabilities (\(\alpha\) = .46--.64; \(\omega\) = .53--.70) fall below conventional benchmarks and limit individual-level interpretation. Three-item scales make high coefficients difficult to attain, but brevity does not make the observed precision adequate. \citet{Loevinger1957Objective} cautioned that highly homogeneous item sets can reduce a construct to repeated versions of a narrow question and argued for broad content representation; that principle supports balancing coverage and homogeneity, not treating low reliability as harmless. The six- and nine-item HPT scores were more precise (\(\omega\) = .66 and .70), so we recommend composite scores rather than individual subscales for group-level research. Associations involving the least reliable measures, especially SDR-5, should be read cautiously because attenuation may obscure weak relationships.

\hypertarget{measurement-fairness-evidence}{%
\subsubsection{5.4 Measurement Fairness Evidence}\label{measurement-fairness-evidence}}

Fit-index changes were compatible with scalar invariance for both
ideology groupings, but inadmissible baseline solutions in the
composite-group analysis prevent a firm conclusion. The DIF analysis
adds item-level evidence, but no design-specific power analysis was
conducted and the unidimensional GRM does not match the CFA structure.
Together, these results provide qualified rather than definitive
fairness evidence.

The correlation pattern---positive scenario-relevant knowledge--HPT
associations \citep{Chapman2021Introduction,Hartmann_Hasselhorn2008Historical,Wineburg2001Historical}, null ideology correlations, null social desirability
correlations---is consistent with the intended score interpretation
\citep{AERA2014Standards,Messick1995Validity}. However, the convergent evidence is bounded by the
scenario-proximal nature of the knowledge measure (see Section 4.9),
which limits the strength of the knowledge-dependent competence claim.

The near-zero ideology--HPT correlation (\emph{r} = $-$.05) should
not be interpreted as conclusive evidence that ideology plays no role
in HPT performance. As discussed in Section 2.3, two opposing
pathways are theoretically plausible: a congruence pathway (whereby
attitudinal alignment inflates scores) and a reduced integrative
complexity pathway (whereby authoritarian cognition depresses
contextualization). If both pathways operate simultaneously, their
effects could cancel, producing a near-zero net correlation that masks
two moderate-sized opposing effects. The present data cannot
distinguish a genuine null from cancelling pathways.

The present fairness evidence applies to mainstream Czech adolescent
samples with restricted ideological range. Users in contexts with
greater political polarization should conduct their own DIF and
invariance analyses before assuming equivalent measurement.

\hypertarget{implications-for-assessment-practice}{%
\subsubsection{5.5 Implications for Assessment Practice}\label{implications-for-assessment-practice}}

The core practical message is that politically charged assessment content does not automatically introduce construct-irrelevant variance---but this conclusion is bounded by the present sample's restricted ideological range and should not be generalized beyond structured, closed-ended formats without further evidence.

The evaluation protocol used here provides an evidence-based model for evaluating fairness when assessment content intersects with respondents' attitudes. This protocol is not specific to HPT or history education; it applies to any scenario-based assessment where content carries attitudinal valence---moral reasoning instruments using politically divisive dilemmas, science literacy assessments touching evolution or climate change, and civic education assessments involving controversial policies. The general approach involves five steps:

\begin{enumerate}
\def\labelenumi{\arabic{enumi}.}
\tightlist
\item
  Identify the attitudinal dimension theoretically most likely to produce CIV.
\item
  Form extreme groups on that dimension.
\item
  Test DIF at the item level.
\item
  Test measurement invariance at the model level.
\item
  Examine a prespecified set of expected relationships, including whether HPT dimensions relate differentially to external outcomes as theory predicts.
\end{enumerate}

The format itself may be a fairness mechanism. If scenario-based formats channel deliberative processing and override automatic attitude-retrieval \citep{Evans2013}, then the instrument format---not just the construct---buffers against motivated reasoning. However, this hypothesis has an important implication: if the format itself suppresses attitudinal responding, then the present fairness conclusion is conditional on the closed-ended Likert format and cannot be extended to open-ended HPT assessments where ideology could operate through different cognitive pathways.

Several design features may explain the instrument's resistance to
attitudinal contamination: the scenario provides contextual
information that channels deliberative processing, the Likert items
require evaluating propositions about a historical figure rather than
expressing personal attitudes, and the school assessment context may
activate an academic frame. These are hypotheses, not
conclusions---the present study was not designed to test them
experimentally, and doing so would require comparing structured and
open-ended formats using the same content.

An alternative explanation for the observed measurement equivalence
deserves consideration. In a school context where the scenario
involves a Nazi-adjacent historical figure, students may uniformly
avoid endorsing presentist responses---not because they are engaging
in genuine historical reasoning, but because contextualization is the
socially safe response regardless of one's actual political
attitudes. If task-demand compliance produces equivalent response
patterns across ideology groups, the resulting measurement invariance
would reflect shared social desirability pressure rather than genuine
absence of ideological influence on reasoning. The null SDR-5
correlations provide partial evidence against this interpretation, but
a brief self-report measure of social desirability may not capture the
specific task demands operating in a classroom assessment of
politically sensitive historical content.

\hypertarget{limitations}{%
\subsubsection{5.6 Limitations}\label{limitations}}

Several limitations qualify these findings.

\begin{enumerate}
\def\labelenumi{\arabic{enumi}.}

\item \emph{Single country and scenario.} Nine ratings of one Weimar
  scenario do not provide nine independent samples of the HPT content
  domain, so scenario-specific wording, knowledge, or moral salience
  cannot be separated from HPT. Future studies should use multiple
  scenarios spanning different periods, places, and political contexts.

\item \emph{Restricted ideological range.} Depending on the measure,
  18.0\% (FR-LF), 37.5\% (KSA-3), or 25.1\% (their raw-score mean) of respondents
  scored above the scale midpoint, so the DIF and
  invariance analyses compare moderately elevated to low ideology, not
  genuine political extremism. The null ideology finding should be
  interpreted as ``no evidence of CIV in this range-restricted sample''
  rather than as a general demonstration that politically charged
  content is fair.

\item \emph{Limited DIF power.} With \emph{n} \(\approx\) 96 per
  group, DIF power was limited. A formal power simulation for the GRM
  likelihood-ratio test at Bonferroni-corrected \(\alpha\) was not
  conducted, so the minimum detectable DIF magnitude is unknown.

\item \emph{Gender imbalance.} The sample was 65.5\% male, and no DIF
  or sensitivity analysis by gender was conducted. The fairness
  conclusion should therefore be replicated in a more gender-balanced
  sample.

\item \emph{Modest reliability.} Subscale reliabilities (\(\alpha\) =
  .47--.64) constrain individual-level score interpretations; future
  adaptations should consider expanding subscale length. Omega
  hierarchical (\(\omega_h\)) from the bifactor model was not
  computed; this index would clarify what proportion of subscale
  variance is attributable to the general HPT factor versus specific
  subscale factors.

\item \emph{Method and format constraints.} All measures were
  self-report instruments administered in a single session. The
  closed-ended Likert format may not generalise to open-ended HPT
  assessments where ideology could operate through different cognitive
  mechanisms \citep{Smith2017Cognitive,Smith_et_al2018History}.

\item \emph{No response process evidence.} Think-aloud or cognitive-
  interview studies, extending \citet{Huijgen_et_al2017Toward}, could
  test whether students use the intended historical reasoning rather
  than attitude-congruent or task-demand shortcuts.

\end{enumerate}

\hypertarget{recommendations}{%
\subsubsection{5.7 Recommendations}\label{recommendations}}

For researchers developing or adapting HPT instruments: (a) report \(\omega\) alongside \(\alpha\) for short subscales, with Spearman-Brown projections for context; (b) evaluate fairness through DIF and measurement invariance for theoretically relevant grouping variables; (c) consider revising weak-loading items (POP2, ROA2) while preserving the theoretical structure; (d) expand subscale length to improve reliability; and (e) replicate fairness evaluations across diverse samples, scenarios, and cultural contexts.

\hypertarget{conclusion}{%
\subsubsection{5.8 Conclusion}\label{conclusion}}

The Czech data were compatible with a correlated three-factor
representation of the Hartmann and Hasselhorn HPT item sets and produced scores
that relate to scenario-relevant historical knowledge in theoretically
expected ways. Fit-index patterns were compatible with measurement
invariance across ideology groups, but model inadmissibility prevents
a firm equivalence conclusion. At the
sensitivity levels afforded by the present sample and in the observed
ideological range, political ideology was unrelated to HPT scores
across all model specifications.

For researchers and assessment developers, the practical implication
is that structured scenario-based assessments can employ politically
sensitive historical content without detectable measurement
bias---at least in mainstream classroom settings with restricted
ideological range. Whether this conclusion extends to more polarized
populations, open-ended assessment formats, or different historical
scenarios remains an open empirical question. The evaluation protocol
illustrated here---CFA, DIF, MG-CFA, and a prespecified network of
external relationships---provides a possible template for answering that
question wherever attitudinally charged content meets competence
assessment.

\hypertarget{data-availability-statement}{%
\subsection{Data Availability Statement}\label{data-availability-statement}}

De-identified data, annotated analysis scripts, the complete instrument battery (in Czech), supplementary materials, and all deviations from the original instrument are available on the Open Science Framework (https://doi.org/10.17605/OSF.IO/YNG37). Data are provided in machine-readable formats (RDS and xlsx) accompanied by a full variable codebook (PDF). The analysis pipeline consists of seven numbered R Markdown scripts that reproduce all reported results; a README file documents the directory structure, execution order, software requirements, and mapping of scripts to manuscript tables and figures. Supplementary materials (Tables S1--S5) are included in the OSF repository. The study design, hypotheses, and analysis plan were preregistered on OSF prior to data collection; deviations and omissions are disclosed in Table S5.

\hypertarget{conflict-of-interest-disclosure}{%
\subsection{Conflict of Interest Disclosure}\label{conflict-of-interest-disclosure}}

The author declares that he complies with the PCI rule of having no financial conflicts of interest in relation to the content of the article, and has no non-financial conflicts of interest.

\hypertarget{ethics-statement}{%
\subsection{Ethics Statement}\label{ethics-statement}}

This study involved secondary school students completing questionnaires during regular class periods. Under Czech law (Act No.~110/2019 Sb., on the processing of personal data), ethics committee approval is not required for non-interventional educational research involving anonymous questionnaires where no personal data are collected. Participation was voluntary; students were informed about the study's purpose and their right to decline participation without consequences. All responses were de-identified at the point of data entry. No personally identifiable information was collected or stored. The study followed the principles of the Declaration of Helsinki for research involving human participants.

\hypertarget{funding}{%
\subsection{Funding}\label{funding}}

This research received no specific grant from any funding agency in the public, commercial, or not-for-profit sectors.

\hypertarget{author-contributions-credit}{%
\subsection{Author Contributions (CRediT)}\label{author-contributions-credit}}

Juda Kaleta: Conceptualization, Methodology, Software, Formal analysis, Investigation, Data curation, Writing --- original draft, Writing --- review \& editing, Visualization, Project administration.

\hypertarget{acknowledgements}{%
\subsection{Acknowledgements}\label{acknowledgements}}

The author thanks the participating schools, teachers, and students for their cooperation. The translation panel members and expert reviewers who contributed to the adaptation process are acknowledged with gratitude.

\hypertarget{ai-use-disclosure}{%
\subsection{AI Use Disclosure}\label{ai-use-disclosure}}

The following AI tools were used during the preparation of this manuscript:

\begin{itemize}
\tightlist
\item
  \emph{Brainstorming and scripting assistance:} Claude (Anthropic) ---
  used for exploring conceptual framing, structural alternatives, and
  scripting support for data processing and analysis workflows.
\item
  \emph{Language editing:} Grammarly --- used for grammar and style
  checking.
\end{itemize}

All research design decisions, analytical choices, theoretical interpretations, and the final text are solely the author's responsibility. The author reviewed, revised, and approved all content.

\bibliographystyle{apacite}
\bibliography{bibliography,extras}

\newpage

\hypertarget{top-checklist}{%
\section*{TOP Checklist and Disclosures Form}\label{top-checklist}}

This checklist must be completed and included as an appendix at the end of your preprint prior to submission to PCI Psychology. Manuscripts without this document included as an appendix will be returned to authors without review.

The policy of PCI Psychology is to recommend papers only if the data, methods used in the analysis, and any digital materials used to conduct the research are clearly and precisely documented and are maximally available to any researcher for purposes of reproducing the results or replicating the procedure. PCI Psych follows the principle of ``as open as possible, as closed as necessary.''

\bigskip

\textbf{First author name (last/family, first/given):} Kaleta, Juda

\textbf{Preprint DOI or URL:} \url{https://doi.org/10.17605/OSF.IO/YNG37}

\bigskip

\hypertarget{section-1-data}{%
\subsection*{Section 1: Data}\label{section-1-data}}

\textbf{Does your manuscript contain reports of any data?}

\textbf{[X] Yes} (continue with next question)

\quad $\square$ No (skip to Section 2)

\bigskip

\textbf{Are appropriately anonymized raw data available within a trusted digital repository?}

\textbf{[X] Yes, available at this link:} \url{https://doi.org/10.17605/OSF.IO/YNG37}

\quad $\square$ No, justification:

\bigskip

\textbf{Are third-party data cited in the manuscript, with a DOI?} (e.g., for preexisting data, data deposited in a repository)

\quad $\square$ Yes, the DOI is as follows:

\textbf{[X] No, justification:} The data were collected specifically for this study; no third-party datasets were used.

\bigskip

\textbf{Is there a data dictionary and/or readme file included with the data to make it interpretable?}

\textbf{[X] Yes, available at this link:} \url{https://doi.org/10.17605/OSF.IO/YNG37} (full variable codebook in PDF and README file)

\quad $\square$ No, justification:

\bigskip

\textbf{Do you indicate in the manuscript how the sample size was determined?}

\textbf{[X] Yes.} (Section 3.4: a priori power analysis targeting 80\% power to detect $\beta \geq .15$)

\quad $\square$ No, justification:

\bigskip

\textbf{Do you report all data exclusions (e.g., outliers, careless responders)?}

\textbf{[X] Yes.} (Section 3.4 and Table S5: the middle ideology tertile was excluded for post-registration DIF and MG-CFA analyses; this analysis-specific exclusion is reported.)

\quad $\square$ No, justification:

\bigskip

\textbf{Do you report all inclusion/exclusion criteria and when they were established?}

\textbf{[X] Yes.} (Participant eligibility is described in Section 3.2 and was preregistered; post-registration analysis-specific exclusions are described in Section 3.4 and Table S5.)

\quad $\square$ No, justification:

\bigskip

\textbf{Are all measures, questions, and/or conditions used in the study described in the manuscript or available in the supplemental material?}

\textbf{[X] Yes.} (All instruments described in Sections 3.1--3.3; full Czech-language instrument battery available on OSF: \url{https://doi.org/10.17605/OSF.IO/YNG37})

\quad $\square$ No, justification:

\bigskip

\hypertarget{section-2-analysis}{%
\subsection*{Section 2: Analysis Scripts/Code/Codebooks}\label{section-2-analysis}}

\textbf{Does your manuscript contain any analysis of quantitative or qualitative data?}

\textbf{[X] Yes} (continue with next question)

\quad $\square$ No (skip to Section 3)

\bigskip

\textbf{Are third-party analysis scripts/code (e.g., R, Stata), codebooks, or other relevant documentation available within a trusted digital repository?}

\textbf{[X] Yes, available at this link:} \url{https://doi.org/10.17605/OSF.IO/YNG37} (7 R Markdown pipeline scripts, 4 figure scripts, 2 supplementary analysis scripts, and a Makefile)

\quad $\square$ No, justification:

\bigskip

\textbf{Are the analysis scripts/code (e.g., R, Stata), codebooks, or other relevant documentation cited in the manuscript, with a DOI?}

\textbf{[X] Yes, the DOI is as follows:} \url{https://doi.org/10.17605/OSF.IO/YNG37}

\quad $\square$ No, justification:

\bigskip

\hypertarget{section-3-materials}{%
\subsection*{Section 3: Study Materials}\label{section-3-materials}}

\textbf{Does your manuscript contain any research materials (e.g., stimuli, programming code, questionnaires, interview protocols)?}

\textbf{[X] Yes} (continue with next question)

\quad $\square$ No (skip to Section 4)

\bigskip

\textbf{Are all study materials and descriptions of study procedures available within a trusted digital repository?}

\textbf{[X] Yes, available at this link:} \url{https://doi.org/10.17605/OSF.IO/YNG37} (Czech-language versions of all questionnaires: HPT, FR-LF, KSA-3, SDR-5, and historical knowledge test)

\quad $\square$ No, justification:

\bigskip

\textbf{Are all third-party study materials, descriptions of study procedures, or other relevant documents cited in the manuscript, with a DOI?}

\quad $\square$ Yes, the DOI is as follows:

\textbf{[X] No, justification:} Third-party instruments are cited via standard bibliography references (Hartmann \& Hasselhorn, 2008; Decker et al., 2013; Beierlein et al., 2014; Hays et al., 1989); DOIs are included where available. The adapted Czech versions are deposited on OSF (\url{https://doi.org/10.17605/OSF.IO/YNG37}).

\bigskip

\hypertarget{section-4-preregistration}{%
\subsection*{Section 4: Preregistration}\label{section-4-preregistration}}

\textbf{Were any aspects of your manuscript preregistered?}

\textbf{[X] Yes} (continue with next question)

\quad $\square$ No (do not complete the rest of the form)

\bigskip

\textbf{Does the manuscript contain an accessible link to the preregistration?}

\textbf{[X] Yes, available at this link:} \url{https://osf.io/zsngy}

\quad $\square$ No, justification:

\bigskip

\textbf{Do you clearly indicate in the manuscript which parts were preregistered and which parts were not?}

\textbf{[X] Yes.} (The six registered hypotheses are summarized in Section 2.5. Table S5 distinguishes registered tests, omissions, implemented modifications, and post-registration analyses.)

\quad $\square$ No, justification:

\bigskip

\textbf{Are all preregistered analyses reported in the text or linked in the supplemental material?}

\quad $\square$ Yes.

\textbf{[X] No, justification:} The registered continuous-ideology MIMIC model, complete moderation and simple-slope analyses, clustered correlation comparison, complete covariate set, one-tailed tests, Benjamini--Hochberg correction, and registered missing-data rule were not all implemented. Table S5 reports these omissions and the analyses used instead.

\bigskip

\textbf{Are all deviations from the preregistration plan clearly disclosed in the manuscript (either in text or in a table)?}

\textbf{[X] Yes.} (Section 2.5 and Table S5 disclose the changed inferential hierarchy, omitted registered procedures, instrument discrepancies, and post-registration additions. TOST and its $\pm$0.20 SESOI are explicitly labelled exploratory.)

\end{document}
